% !TeX program = xelatex
% 刘丹丹 — 中文学术简历（与英文版同步）
% Adapted from Pascal Michaillat's minimalist academic CV template.
% https://github.com/pmichaillat/latex-cv
% https://pascalmichaillat.org/f/
% Compile twice with: xelatex main.tex
% Self-contained source; uses standard Fandol, TeX Gyre Pagella, and Latin Modern fonts.
% Content snapshot: 12 November 2026, from the website's about/CV/publication records.
% Publication author order and abbreviated author lists are retained from those records.
% Omitted: inconsistent master's thesis date (title retained).
% Template license follows.
% MIT License
%
% Copyright (c) 2025–present Pascal Michaillat
%
% Permission is hereby granted, free of charge, to any person obtaining a copy
% of this software and associated documentation files (the "Software"), to deal
% in the Software without restriction, including without limitation the rights
% to use, copy, modify, merge, publish, distribute, sublicense, and/or sell
% copies of the Software, and to permit persons to whom the Software is
% furnished to do so, subject to the following conditions:
%
% The above copyright notice and this permission notice shall be included in all
% copies or substantial portions of the Software.
%
% THE SOFTWARE IS PROVIDED "AS IS", WITHOUT WARRANTY OF ANY KIND, EXPRESS OR
% IMPLIED, INCLUDING BUT NOT LIMITED TO THE WARRANTIES OF MERCHANTABILITY,
% FITNESS FOR A PARTICULAR PURPOSE AND NONINFRINGEMENT. IN NO EVENT SHALL THE
% AUTHORS OR COPYRIGHT HOLDERS BE LIABLE FOR ANY CLAIM, DAMAGES OR OTHER
% LIABILITY, WHETHER IN AN ACTION OF CONTRACT, TORT OR OTHERWISE, ARISING FROM,
% OUT OF OR IN CONNECTION WITH THE SOFTWARE OR THE USE OR OTHER DEALINGS IN THE
% SOFTWARE.
%

\documentclass[a4paper,11pt]{article}
\usepackage[margin=1.9cm,top=1.65cm,bottom=1.85cm,headheight=14pt,headsep=12pt,footskip=24pt]{geometry}
\usepackage{fontspec,xeCJK,microtype,xcolor,fancyhdr,titlesec,needspace}
\setCJKmainfont{FandolSong-Regular.otf}[BoldFont=FandolSong-Bold.otf,ItalicFont=FandolKai-Regular.otf]
\setCJKsansfont{FandolHei-Regular.otf}[BoldFont=FandolHei-Bold.otf]
\usepackage[hidelinks,pdfdisplaydoctitle=true]{hyperref}
\setmainfont{texgyrepagella-regular.otf}[BoldFont=texgyrepagella-bold.otf,ItalicFont=texgyrepagella-italic.otf,BoldItalicFont=texgyrepagella-bolditalic.otf]
\setsansfont{lmsans10-regular.otf}[BoldFont=lmsans10-bold.otf,ItalicFont=lmsans10-oblique.otf,BoldItalicFont=lmsans10-boldoblique.otf]
\definecolor{navy}{HTML}{173E52}
\definecolor{teal}{HTML}{367B88}
\definecolor{ink}{HTML}{242D32}
\definecolor{muted}{HTML}{54636C}
\color{ink}
\setlength{\parindent}{0pt}
\setlength{\parskip}{5pt}
\linespread{1.15}
\AtBeginDocument{\frenchspacing}
\hypersetup{pdftitle={刘丹丹 — 中文学术简历},pdfauthor={刘丹丹 / Dandan Liu},
  pdfsubject={研究方向、教育背景与学术论文}}
\titleformat*{\subsection}{\large\sffamily\bfseries\color{navy}}
\titlespacing*{\subsection}{0pt}{4pt}{5pt}
\newcommand{\sep}{{\,\textperiodcentered\,\,}}
\newcommand{\hsep}{\noindent\textcolor{teal}{\rule{\textwidth}{0.5pt}}}
\newcommand{\block}[1]{\par\addvspace{7pt}\Needspace{5\baselineskip}\hsep\vspace{-8pt}\subsection*{#1}}
\newcommand{\cvfoot}[2]{\fancyfoot[#1]{\footnotesize\textcolor{muted}{#2}}}
\pagestyle{fancy}
\fancyhf{}
\renewcommand{\headrulewidth}{0pt}
\fancyhead[L]{\footnotesize\textcolor{muted}{刘丹丹 / Dandan Liu}}
\fancyhead[R]{\footnotesize\textcolor{muted}{研究与学术成果}}
\cvfoot{L}{\href{mailto:todandan@foxmail.com}{todandan@foxmail.com}}
\cvfoot{C}{2026年11月12日}
\cvfoot{R}{\thepage\,/\,\pageref*{LastPage}}
\fancypagestyle{opening}{\fancyhead{}}
\newcommand{\publication}[2]{%
  \par\addvspace{8pt}\noindent
  \begin{minipage}{\linewidth}
    \linespread{1.06}\fontsize{10.8}{13.4}\selectfont\raggedright
    \hangindent=2.4em\hangafter=1
    \makebox[2.4em][l]{\textcolor{teal}{[#1]}}#2
  \end{minipage}\par}
\begin{document}
\thispagestyle{opening}
{\sffamily\small\color{muted} 学术简历 · 中文版}\par
\vspace{1mm}
{\fontsize{30}{36}\selectfont\color{navy}\textbf{刘丹丹}}\quad
{\fontsize{18}{24}\selectfont\color{navy}Dandan Liu}\par
{\large 媒体与传播学博士研究生}\par
{\small 马来亚大学 · 文学与社会科学学院 · 媒体与传播研究系\\
马来西亚，吉隆坡}\par
{\small\href{mailto:todandan@foxmail.com}{todandan@foxmail.com} \sep
\href{https://s2134717ldd.github.io/smiledan713.github.io/}{学术主页} \sep
\href{https://orcid.org/0009-0002-4702-8070}{ORCID: 0009-0002-4702-8070}}\par

\vspace{3mm}
\begingroup
\setlength{\fboxsep}{12pt}
\noindent\colorbox{navy}{%
\begin{minipage}{\dimexpr\linewidth-2\fboxsep\relax}
\color{white}
{\sffamily\fontsize{10}{12}\selectfont 研究兴趣}\par\vspace{4pt}
{\fontsize{14}{21}\selectfont
计算传播学 \sep 社交媒体 \sep 危机传播\\
人机交互 \sep 数字话语与媒介动态}
\end{minipage}}
\endgroup

\block{研究概述}
我关注数字平台如何塑造公共话语与人际互动，结合质性研究与计算方法，研究社交媒体、危机传播和人机交互。

\block{研究方向}
\textbf{数字话语与危机传播：}关注疫情传播与跨国因素的影响，考察政府信息、社交媒体与公共话语的关系。

\textbf{人机传播：}研究公共社交平台中的大语言模型智能体、情感支持与真实性，以及人际对话结构的变化。

\textbf{平台治理与信任：}关注公平性披露、内容溯源、用户选择，以及人工智能生成摘要和可穿戴设备健康评分的权威性。

\block{教育背景}
\textbf{媒体与传播学博士研究生}\hfill 2021年--2027年5月（预计）\\
马来亚大学，文学与社会科学学院，媒体与传播研究系。\\
{\small 博士论文：《中国新冠疫情新闻发布会中的语义治理与危机框架（2020--2023年）》\\
英文题目：\emph{SEMANTIC GOVERNANCE AND CRISIS FRAMING IN CHINA'S COVID-19 PRESS CONFERENCES 2020--2023}.}

\textbf{新闻与传播学硕士}\hfill 2016--2018年\\
复旦大学，新闻学院。\\
{\small 硕士论文：《“罗辑思维”的媒体属性与社群属性》。}

\textbf{新闻学学士}\hfill 2010--2014年\\
复旦大学，新闻学院。\\
{\small 本科论文：《中国大学生手机新闻客户端使用与满足研究》。}

\clearpage
\block{第一作者论文（11篇）}
{\small\color{muted}作者顺序、论文题名与发表信息保留原文；粗体标注本人。}\par
% Source record: _publications/2026-09-18-ubicomp-wearable-wellness-scores.md
\publication{1}{\textbf{Liu, D.}, \& Fan, G. (2027). When the Score Speaks First: Negotiating the Authority of Wearable Wellness Scores in Everyday Life. \emph{ACM International Joint Conference on Pervasive and Ubiquitous Computing (UbiComp/ISWC 2027)}. 已接收，待刊。}

% Source record: _publications/2026-09-14-cscw-continuity-work.md
\publication{2}{\textbf{Liu, D.}, Sabri, A.Q.M., Rahman, S.S.A., Pan, L., \& Fan, G. (2027). When AI Breaks Down: Continuity Work, Dependency, and Repair in AI-Mediated Collaboration. \emph{ACM Conference on Computer-Supported Cooperative Work and Social Computing (CSCW 2027)}. 已接收，待刊。}

% Source record: _publications/2026-08-22-shared-space.md
\publication{3}{\textbf{Liu, D.}, Pan, L., Md Sabri, A.Q., \& Fan, G. (2026). From Solo Post to Shared Space: How a Public LLM Agent Reshapes Human-to-Human Conversation Structure on a Social Platform. \emph{ACM Transactions on Social Computing}. \href{https://doi.org/10.1145/3844666}{doi:10.1145/3844666}.}

% Source record: _publications/2026-07-15-fairness-labels.md
\publication{4}{\textbf{Liu, D.}, \& Fan, G. (2026). ``Fairness Labels'' in Feeds: Process Disclosures, Provenance, and User Choice Among Short-Video Creators in Mainland China. \emph{Social Media + Society}, 12(3), Article 20563051261475922. \href{https://doi.org/10.1177/20563051261475922}{doi:10.1177/20563051261475922}.}

% Source record: _publications/2026-04-03-textless-tiny-private.md
\publication{5}{\textbf{Liu, D.}, Fan, G., Guo, Q., Zhang, R., \& Pan, L. (2026). Textless, Tiny, and Private: Hashed n-Grams with LLM-Teacher Distillation and Local Differential Privacy for Mental-Health Screening on Social Media. \emph{IEEE Transactions on Computational Social Systems} （在线优先发表）. \href{https://doi.org/10.1109/TCSS.2026.3693692}{doi:10.1109/TCSS.2026.3693692}.}

% Source record: _publications/2026-01-01-tool-tutor-crutch.md
\publication{6}{\textbf{Liu, D.}, Fan, G., \& Pan, L. (2026). Tool, tutor, or crutch?: A grounded theory of cognitive scaffolding and offloading in AI-assisted programming education. \emph{International Journal of STEM Education}, 13, 10. \href{https://doi.org/10.1186/s40594-025-00592-w}{doi:10.1186/s40594-025-00592-w}.}

% Source record: _publications/2025-intrac-frame-dynamics.md
\publication{7}{\textbf{Liu, DanDan}. (2025). Beyond Top-Down Messaging: Frame Dynamics Across Government, Social Media, and Public Discourse in China's COVID-19 Crisis Communication. In \emph{INTRAC 2025}.}

% Source record: _publications/2025-zero-covid-china.md
\publication{8}{\textbf{Liu, Dandan}, Siti Ezaleila Binti Mustafa, and Nor Zaliza Binti Sarmiti. (2025). From zero-COVID to global alignment: transnational pressures and China's pandemic communication transformation. \emph{Globalization and Health}. \href{https://doi.org/10.1186/s12992-025-01155-y}{doi:10.1186/s12992-025-01155-y}.}

% Source record: _publications/2025-cogsci-ai-voices.md
\publication{9}{\textbf{Liu, D.}, et al. (2025). When Machines Speak with Feeling: Investigating Emotional Prosody, Authenticity, and Trust in AI vs. Human Voices. In \emph{Proceedings of the Annual Meeting of the Cognitive Science Society (Vol. 47)}.}

% Source record: _publications/2025-cogsci-llm-emotional-support.md
\publication{10}{\textbf{Liu, D.}, et al. (2025). Co-Constructing Meaning with Large Language Models: A Longitudinal Analysis of Human-AI Dialogues in Emotional Support Contexts. In \emph{Proceedings of the Annual Meeting of the Cognitive Science Society (Vol. 47)}.}

% Source record: _publications/2025-04-24-digital-duality.md
\publication{11}{\textbf{Liu, D.}, \& Fan, G. (2025). Digital duality: negotiating tradition, modernity, and identity among rural housewife Douyin creators in China. \emph{Information, Communication and Society}. \href{https://doi.org/10.1080/1369118x.2025.2496903}{doi:10.1080/1369118x.2025.2496903}.}

\clearpage
\block{合作论文：2026年}
% Source record: _publications/2026-08-20-fairsynth.md
\publication{1}{Fan, G., \textbf{Liu, D.}, Guo, Q., Niu, X., \& Pan, L. (2026). FairSynth: Instrument-Relative Evaluation and Calibration of LLM Survey Simulations for Measurement Fairness. In \emph{Findings of the Association for Computational Linguistics: EMNLP 2026}. 已接收。}

% Source record: _publications/2026-08-01-feeling-rules.md
\publication{2}{Fan, G., \textbf{Liu, D.}, Sabri, A.Q.M., Zhang, R., \& Pan, L. (2026). Feeling Rules in Language Models: Mapping Norms of Emotional Appropriateness Across Roles, Institutions, and Intensity. In \emph{Proceedings of the 64th Annual Meeting of the Association for Computational Linguistics (ACL 2026)}. ACL.}

% Source record: _publications/2026-07-06-graph-preference-learning.md
\publication{3}{Fan, G., \textbf{Liu, D.}, Sabri, A.Q.M., \& Pan, L. (2026). Graph-Preference Learning: Debiasing Network-Sampled Human Feedback for Target Welfare Estimation. In \emph{Proceedings of the 43rd International Conference on Machine Learning (ICML 2026)}. PMLR.}

% Source record: _publications/2026-07-01-multi-llm-persona.md
\publication{4}{Fan, G., \textbf{Liu, D.}, Pan, L., Zhang, R., \& Guo, Q. (2026). Multi-LLM Persona Generation for Virtual Focus Groups in Software Engineering: A Controlled, Multi-Domain Study of Emotional Requirements Elicitation. In \emph{Proceedings of the ACM International Conference on the Foundations of Software Engineering (FSE 2026)}. ACM.}

% Source record: _publications/2026-06-22-conformal-k.md
\publication{5}{Guangrui Fan, \textbf{DanDan Liu}, and Lihu Pan. 2026. Conformal@K: Distribution-Free Top-K Miss-Risk Control for Recommendation with Overlapping-Group and Two-Stage Guarantees. \emph{ACM Transactions on Information Systems} 44, 6, Article 138 (July 2026), 42 pages. \href{https://doi.org/10.1145/3828550}{doi:10.1145/3828550}.}

% Source record: _publications/2026-05-03-eco-nudging.md
\publication{6}{Fan, G., \textbf{Liu, D.}, et al. (2026). Co-Adaptive Eco-Nudging: A Privacy-Preserving Contextual Bandit with User-Taught Preferences. In \emph{Proceedings of the 2026 CHI Conference on Human Factors in Computing Systems (CHI 2026)}. ACM.}

% Source record: _publications/2026-05-02-verification-load.md
\publication{7}{Fan, G., \textbf{Liu, D.}, et al. (2026). When Help Hurts: Verification Load and Fatigue with AI Coding Assistants. In \emph{Proceedings of the 2026 CHI Conference on Human Factors in Computing Systems (CHI 2026)}. ACM.}

% Source record: _publications/2026-05-01-ai-authorship-authenticity.md
\publication{8}{Fan, G., \textbf{Liu, D.}, et al. (2026). Is It Still You? Attributing Authorship and Authenticity in AI-Assisted Romantic Communication. In \emph{Proceedings of the 2026 CHI Conference on Human Factors in Computing Systems (CHI 2026)}. ACM.}

% Source record: _publications/2026-04-25-consent-boundaries.md
\publication{9}{Fan, G., \textbf{Liu, D.}, Pan, L., Zhang, R., \& Sabri, A.Q.M. (2026). Consent Boundaries by Play: A CI-Grounded Micro-Game for Eliciting Stable Preferences and Diagnosing Opt-Out Friction in Web Consent. In \emph{Proceedings of the 2026 European Conference on Computer-Supported Cooperative Work (ECSCW 2026)}. EUSSET.}

% Source record: _publications/2026-04-01-audit-of-audits.md
\publication{10}{Fan, G., \textbf{Liu, D.}, Sabri, A.Q.M., \& Pan, L. (2026). Audit-of-Audits for the Web: Bayesian Meta-Evaluation that Yields Interval-Valued, Threshold-Aligned Fairness Claims. In \emph{Proceedings of The Web Conference 2026 (WWW 2026)}.}

% Source record: _publications/2026-02-01-mind-the-gap.md
\publication{11}{Fan, G., \textbf{Liu, D.}, \& Pan, L. (2026). Mind the Gap: Predicting, Explaining and Reducing Time-to-First-Comment (Reply Gap) in Online Mental-Health Communities. In \emph{Proceedings of the 40th AAAI Conference on Artificial Intelligence (AAAI 2026)}.}

\clearpage
\block{合作论文：2025年}
% Source record: _publications/2025-ijcai-creative-momentum.md
\publication{12}{Fan, G., \textbf{Liu, D.}, Pan, L., \& Huang, Y. (2025). Creative Momentum Transfer: How Timing and Labeling of AI Suggestions Shape Iterative Human Ideation. In \emph{International Joint Conferences on Artificial Intelligence (IJCAI)}.}

% Source record: _publications/2025-skim-or-swim.md
\publication{13}{Fan, G., \textbf{Liu, D.}, \& Huang, Y. (2025). Skim or Swim? Investigating AI-Generated Summaries, Trust, and Comprehension in Chinese Mobile Reading. \emph{International Journal of Human-Computer Interaction}, 1-22. \href{https://doi.org/10.1080/10447318.2025.2574513}{doi:10.1080/10447318.2025.2574513}.}

% Source record: _publications/2025-ijcnn-artists-ai.md
\publication{14}{Fan, G., \textbf{Liu, D.}, Pan, L., \& Huang, Y. (2025). Mapping Override Behavior: Investigating Why and How Artists Reject AI Suggestions in Collaborative Creation. In \emph{2025 International Joint Conference on Neural Networks (IJCNN)}.}

% Source record: _publications/2025-03-04-ai-pair-programming.md
\publication{15}{Fan, G., \textbf{Liu, D.}, Zhang, R., \& Pan, L. (2025). The impact of AI-assisted pair programming on student motivation, programming anxiety, collaborative learning, and programming performance: a comparative study with traditional pair programming and individual approaches. \emph{International Journal of STEM Education, 12}(1), 16. \href{https://doi.org/10.1186/s40594-025-00537-3}{doi:10.1186/s40594-025-00537-3}.}

% Source record: _publications/2025-03-01-emotional-labor-nurses.md
\publication{16}{Yang, Fang, \textbf{Dandan Liu}, and Guangrui Fan. (2025). Emotional labor and coping strategies of gynecological nurses in recurrent pregnancy loss care: a qualitative phenomenological study. \emph{BMC Nursing}, 24(1), 1-11. \href{https://doi.org/10.1186/s12912-025-02884-6}{doi:10.1186/s12912-025-02884-6}.}

% Source record: _publications/2025-01-11-power-posing.md
\publication{17}{Deng, W., \textbf{Liu, D.}, \& Fan, G. (2025). The impact of power posing on presentation performance and self-efficacy in higher education: An embodied cognition approach. \emph{Innovations in Education and Teaching International}. \href{https://doi.org/10.1080/14703297.2025.2452393}{doi:10.1080/14703297.2025.2452393}.}

% Source record: _publications/2025-wechat-prenatal-education.md
\publication{18}{He, Y., Fan, G., Fan, G., \& \textbf{Liu, D.} (2025). Exploring nurse and patient perspectives on WeChat-based prenatal education in Chinese public hospitals: a qualitative inquiry. \emph{BMC Nursing}, 24(1), 459.}

\label{LastPage}
\end{document}

